\documentclass[journal]{IEEEtran}
\pdfoutput=1


\usepackage{amsmath,amssymb,amsfonts}
\usepackage{algorithmic}
\usepackage{graphicx}
\usepackage{textcomp}
\usepackage{xcolor}
\usepackage{booktabs}
\usepackage{multirow}
\usepackage{cite}
\usepackage{url}
\usepackage{hyperref}
\usepackage{bm}
\usepackage{microtype}

\hypersetup{colorlinks=true,linkcolor=blue,citecolor=blue,urlcolor=blue}

% Spacing parameters in preamble so they are not overridden by IEEEtran init
\tolerance=1000
\hbadness=10000
\emergencystretch=1.5em

\begin{document}
\hyphenation{SDFormer ATSR AGF RevIN}

\title{SDFormer: Multi-Granularity Temporal Modeling for Supply-Demand Forecasting}

% \author{Liying~Wang, Peng~Zhang, Cheng~Wang\textsuperscript{*}
% \thanks{The authors are with the School of Computer Science and Technology,
%   Tongji University, Shanghai, China.}
% \thanks{\textsuperscript{*}Corresponding author: wangcheng@tongji.edu.cn}}

\maketitle

% -------------------------------------------------------
\begin{abstract}
Supply-demand sequences in computing power networks (CPNs) are characterized by
the superposition of stationary periodic trends and non-stationary burst
spikes, a mixture that challenges both temporal encoding and input
normalization in existing forecasting methods.
We propose SDFormer, which addresses this through multi-granularity
behavioral decoupling. Parallel patch encoders at fine, medium, and coarse
scales separate spike-dominated and trend-dominated components, while
spike-aware recalibration amplifies burst signals before patch projection,
an adaptive gating mechanism routes each input instance to its most suitable
granularity, and robust median-based normalization prevents burst events
from corrupting the learned representations.
On high-burstiness benchmarks, SDFormer reduces the 96-step MSE by 7.6\%
over PatchTST and 7.5\% over iTransformer, and the 192-step MSE by 14.9\%
over PatchTST, while maintaining competitive performance on smooth datasets.
Code is available at \url{https://github.com/Wingspeg/sdformer}.
\end{abstract}

\begin{IEEEkeywords}
Time-series forecasting, CPNs, multi-granularity
modeling, supply-demand prediction, anomaly-aware normalization, Transformer
\end{IEEEkeywords}

% -------------------------------------------------------
\section{Introduction}

\IEEEPARstart{W}{ith} the rapid expansion of CPNs,
resource supply-demand forecasting has become a critical technology for
capacity planning and elastic scheduling~\cite{zhou2023survey}.
The proliferation of large-scale AI workloads, video streaming, and edge
computing services has further intensified the volatility of demand signals,
making reliable forecasting an increasingly difficult open problem~\cite{alieldin2014burstiness}.

Supply-demand sequences in CPNs exhibit a distinctive
characteristic, namely the temporal superposition of stationary and non-stationary
behavior.
Stationary behavior is driven by business regularity and exhibits stable
diurnal and weekly cycles as low-frequency oscillations whose periodic
structure requires long historical windows to estimate accurately~\cite{moreno2014workload}.
For instance, GPU cluster utilization in a typical CPN follows a predictable
daily rhythm tied to business hours, with a baseline that can be reliably
modeled using periodic decomposition over a multi-day look-back window.
Non-stationary behavior is triggered by stochastic events such as sudden
traffic bursts, batch task submissions, and failure-induced traffic
migrations~\cite{alieldin2014burstiness}.
It manifests as sparse, steep pulse spikes that are short in duration but large in amplitude.
Unlike stationary periodicity, these spikes cannot be anticipated from
historical patterns alone; their onset time, duration, and magnitude are
all highly variable.
The two components are fully mixed in the time domain and cannot be
separated by simple thresholding, since a spike superimposed on a periodic
background is indistinguishable from a high-amplitude period without
additional contextual information.
This creates a dual challenge in which long-range memory capacity and local
detail capture are mutually constrained under any single fixed-window
architecture, while non-stationary spikes contaminate normalization
statistics and compress the feature representation space for stationary
components.

Existing forecasting methods fall short in handling this mixture from
three perspectives.

\textbf{Single-granularity temporal encoding.}
Transformer-based methods such as Informer~\cite{informer},
Autoformer~\cite{autoformer}, and FEDformer~\cite{fedformer} apply
attention at the individual time-step level, lacking the receptive field
diversity needed to simultaneously model slow periodic trends and fast
burst spikes.
PatchTST~\cite{patchtst} improves local semantic capture through patch
tokenization, but a single fixed patch size remains unable to cover both
frequency components at once: a large patch size that captures weekly
periodicity will average over a spike that lasts only a few time steps,
while a small patch size that resolves spikes will fail to accumulate
the long-range context needed for trend modeling.
iTransformer~\cite{itransformer} redirects attention to the variable
dimension, which benefits multivariate correlation modeling but does not
address within-sequence multi-frequency dynamics.
In all of these architectures, the model is forced to make a single
receptive field commitment that inevitably sacrifices either spike
resolution or trend coverage.

\textbf{Inadequate sequence decomposition.}
Linear and MLP-based methods such as DLinear~\cite{dlinear} and
N-HiTS~\cite{nhits} decompose the input into a trend and a remainder
using a single moving-average split.
This two-component representation cannot capture the sparse, steep spikes
that characterize bursty CPN workloads, as such events are absorbed into
the remainder component without dedicated modeling capacity.
The remainder is then modeled by the same linear or MLP layer as
ordinary residual fluctuations, which systematically underestimates
spike peaks because the model has no mechanism to distinguish a spike
from ordinary noise of similar magnitude.
TimeMixer~\cite{timemixer} introduces multi-scale mixing but applies
uniform operations across all time steps, making no distinction between
stationary background steps and non-stationary burst events.
Multi-scale CNN methods such as MICN~\cite{micn} and ModernTCN~\cite{moderntcn}
similarly treat all time steps equally within their convolutional kernels,
providing no mechanism to selectively amplify sparse spike signals.

\textbf{Fragility of instance normalization under burstiness.}
RevIN~\cite{revin} and its variants standardize each input instance by
its sample mean and variance.
Under bursty CPN sequences, sparse extreme values inflate both statistics,
shifting the normalization center and scale away from the values that
describe normal behavior.
Non-stationary Transformer~\cite{nonstationary} identifies
over-stationarization as a failure mode and proposes De-stationary
Attention to restore non-stationary expressiveness, yet it does not
address the corrupted normalization statistics caused by spike
contamination at the input stage.

To address these limitations, this paper proposes \textbf{SDFormer}, which
redesigns the forecasting architecture from a behavioral-decoupling
perspective.
The main contributions are as follows.

\begin{itemize}
  \item We identify that the superposition of stationary trends and
    non-stationary burst spikes in CPN supply-demand sequences poses
    fundamental challenges to both temporal encoding and input
    normalization in existing forecasting methods.

  \item We propose SDFormer, a multi-granularity Transformer that decouples
    behavioral components through parallel independent encoders at fine,
    medium, and coarse scales, spike-aware recalibration, adaptive gating
    fusion, and robust median-based normalization.

  \item Experiments on five benchmarks demonstrate that SDFormer achieves
    state-of-the-art performance on high-burstiness sequences and maintains
    competitive results on smooth datasets.
\end{itemize}

The remainder of this paper is organized as follows.
Section~\ref{sec:related} reviews related work.
Section~\ref{sec:method} presents the SDFormer architecture.
Section~\ref{sec:experiments} reports experimental results.
Section~\ref{sec:conclusion} concludes.

% -------------------------------------------------------
\section{Related Work}
\label{sec:related}

\subsection{Transformer-Based Time Series Forecasting}

The Transformer architecture has been widely explored for long-term time
series forecasting.
Early efforts focused on reducing the quadratic complexity of self-attention:
Informer~\cite{informer} introduced ProbSparse attention, Reformer~\cite{reformer}
adopted locality-sensitive hashing, and Flowformer~\cite{flowformer} linearized
attention via conservation flows.
Subsequent work shifted toward richer temporal representations.
Autoformer~\cite{autoformer} incorporated Auto-Correlation with progressive
decomposition to model seasonality, and FEDformer~\cite{fedformer} extended
this to the frequency domain.
PatchTST~\cite{patchtst} marked a significant step forward by tokenizing
overlapping patches rather than individual time steps, capturing local
semantics and enabling longer look-back windows.
However, a single fixed patch size cannot simultaneously resolve both
low-frequency periodic trends and high-frequency burst spikes, as larger
patches suppress transient events while smaller patches lose long-range
context.
Crossformer~\cite{crossformer} further introduced two-stage attention to
jointly model temporal and cross-variable dependencies.
iTransformer~\cite{itransformer} inverted the attention axis, treating each
variable's full sequence as a token to focus on multivariate correlations.
Non-stationary Transformer~\cite{nonstationary} identified that excessive
normalization destroys non-stationary expressiveness, proposing De-stationary
Attention as a remedy.
Despite these advances, all methods above rely on a single temporal
granularity and do not distinguish between stationary trends and
non-stationary burst spikes within the same window, which limits their
effectiveness on bursty CPN supply-demand sequences.
Their attention mechanisms are optimized for either long-range dependencies
or local patterns, but not both simultaneously.
MSformer~\cite{msformer} specifically targets CPN load forecasting with
multi-scale feature fusion, but fuses features in a shared space without
isolating the feature representations of different frequency components.
SDFormer differs fundamentally by maintaining fully independent encoders
per granularity, eliminating cross-frequency gradient interference during
training.

\subsection{Non-Transformer Forecasting Methods}

Beyond Transformers, linear, MLP, CNN, and recurrent architectures have
demonstrated competitive or superior performance on smooth benchmarks.
DLinear~\cite{dlinear} showed that a simple moving-average decomposition
followed by linear projection can match complex Transformer models on many
standard datasets, calling into question the necessity of attention for
stationary sequences.
N-BEATS~\cite{nbeats} and N-HiTS~\cite{nhits} used fully connected residual
stacks with hierarchical interpolation for interpretable trend-seasonality
decomposition.
TimeMixer~\cite{timemixer} proposed multi-scale mixing with
Past-Decomposable Mixing and Future-Multipredictor Mixing modules.
On the convolution front, ModernTCN~\cite{moderntcn} revived temporal
convolution~\cite{tcn} with large kernels, while MICN~\cite{micn}
combined multi-scale isometric convolutions with local-global fusion.
For sequential modeling, LSTM~\cite{lstm} and DeepAR~\cite{deepar}
established strong recurrent baselines, while Mamba~\cite{mamba} and
S-Mamba~\cite{smamba} introduced selective state space models with
near-linear complexity.
WPMixer~\cite{xpatch} applied wavelet decomposition for multi-resolution
patch mixing.
A common limitation across all these methods is that they treat all time
steps uniformly, without targeted mechanisms to amplify or isolate the
non-stationary burst events that dominate CPN supply-demand dynamics.
Recurrent methods such as LSTM and Mamba process time steps sequentially
with a shared hidden state, so spike information is gradually overwritten
by subsequent smooth steps before reaching the output layer.
Linear methods apply the same weight matrix to all time steps regardless
of their local amplitude, which systematically underweights sparse high-amplitude
events relative to the dense low-amplitude background.
Wavelet-based methods such as WPMixer distribute spike energy across multiple
frequency sub-bands, which improves multi-scale representation but does not
provide dedicated capacity for the spatial localization of individual spike events.
SDFormer addresses this gap through ATSR, which explicitly assigns learnable
per-step importance weights within each patch, and through independent
fine-grained encoders that maintain dedicated capacity for high-frequency
transient modeling without sacrificing coarse-grained trend coverage.

\subsection{Multi-Scale Temporal Modeling}

Capturing dependencies at multiple temporal scales is a recognized
challenge in time series forecasting.
Pyraformer~\cite{pyraformer} addressed this with pyramidal attention at
different resolutions, and TimesNet~\cite{timesnet} reshaped sequences into
two-dimensional structures to exploit multi-periodicity.
Koopa~\cite{koopa} applied Koopman theory to disentangle time-varying and
time-invariant components.
FITS~\cite{fits} showed that frequency interpolation with minimal parameters
suffices for smooth periodic sequences.
However, these methods are motivated by multi-periodicity rather than
behavioral decoupling: they fuse or transform features across scales within
a shared representation space, which means that non-stationary spike signals
remain susceptible to dilution by surrounding stationary context.
In Pyraformer, coarser pyramid levels summarize longer segments by pooling,
which averages over spikes and reduces their amplitude in the summary
representation.
In TimesNet, the 2D reshape is designed to exploit inter-period repetition,
which is advantageous for strictly periodic data but provides no mechanism
for localizing aperiodic spike events within a period.
SDFormer departs from this by using fully independent encoders per
granularity, ensuring that the feature spaces for different frequency
components are completely isolated.
Crucially, the independence is maintained not only in the forward pass but
also during backpropagation: gradients from the spike detection loss do not
flow into the trend modeling parameters and vice versa, allowing each
encoder to converge to a solution that is globally optimal for its
frequency band rather than a compromise between conflicting objectives.

\subsection{Normalization in Time Series Forecasting}

Distribution shift between training and test sequences is a fundamental
challenge in time series forecasting.
RevIN~\cite{revin} addressed this with reversible instance normalization
using per-instance mean and variance, and has since become a standard
component in many architectures.
SAN~\cite{san} extended it to a slice-aware paradigm for local distribution
adaptation.
Non-stationary Transformer~\cite{nonstationary} showed that aggressive
normalization suppresses non-stationary information, motivating its
De-stationary Attention.
A shared weakness of all existing schemes is their reliance on mean and
variance statistics, which are highly sensitive to sparse extreme values.
In bursty CPN sequences, even a small number of spike time steps can
significantly bias these statistics, distorting the normalized feature
space for the remaining stationary steps.
For example, a single spike of amplitude $5\sigma$ in a window of 512
steps shifts the sample mean by approximately $0.01\sigma$ and inflates
the sample variance by roughly $0.05\sigma^2$, which may seem small in
absolute terms but is sufficient to misplace the normalization center for
the stationary background by a perceptible fraction of its dynamic range.
This effect compounds across multiple spikes, and is particularly
problematic for CPN workloads where burst events can be orders of
magnitude larger than baseline utilization.
SDFormer addresses this with Anomaly-Aware Normalization, which replaces
mean and variance with median and MAD estimators that are provably
resistant to up to 50\% spike contamination, and adaptively interpolates
between robust and classical normalization based on the measured kurtosis
and tail probability of each input instance.

% -------------------------------------------------------
\section{Methodology}
\label{sec:method}

Fig.~\ref{fig:arch} shows the overall SDFormer architecture.
Each input sequence passes through Anomaly-Aware Normalization, then enters
three independent patch encoders at fine, medium, and coarse granularities.
ATSR modulates per-step importance within each patch before projection.
AGF fuses the three encoded representations with per-instance gating weights.
A linear projection maps the fused output to the prediction horizon, and
Anomaly-Aware Normalization is inverted to restore the original scale.

\begin{figure}[!t]
  \centering
  \includegraphics[width=\columnwidth]{figures/sdformer_architecture.png}
  \caption{Overall architecture of SDFormer.}
  \label{fig:arch}
\end{figure}

\subsection{Problem Formulation}

Given a supply-demand sequence $\mathbf{X} \in \mathbb{R}^{L \times N}$,
where $L$ is the look-back window length and $N$ is the number of resource
channels, the goal is to predict $\hat{\mathbf{Y}} \in \mathbb{R}^{T \times N}$.
Each input channel is modeled independently following the Channel Independent
principle~\cite{patchtst}, which avoids spurious cross-channel correlations
that can arise when variable-dimension attention is applied to heterogeneous
resource metrics with different units and dynamic ranges.

We characterize each channel as a mixture of two overlapping behavioral components:
\begin{equation}
x_{t,j} = x^{(\mathrm{s})}_{t,j} + x^{(\mathrm{ns})}_{t,j}
\label{eq:decomp}
\end{equation}
where $x^{(\mathrm{s})}_{t,j}$ is the stationary low-frequency component
driven by business periodicity, and $x^{(\mathrm{ns})}_{t,j}$ is the
non-stationary sparse spike component triggered by stochastic events.
Unlike decomposition-based methods that require an explicit separation
stage~\cite{autoformer,dlinear}, SDFormer implicitly allocates modeling
responsibility for each component to a dedicated granularity encoder:
the fine-grained encoder is sensitive to the high-amplitude transients of
$x^{(\mathrm{ns})}_{t,j}$, while the coarse-grained encoder integrates the
extended context needed to model $x^{(\mathrm{s})}_{t,j}$.
This implicit allocation eliminates the decomposition error that propagates
through explicit separation pipelines when the two components partially
overlap in the frequency domain.

\subsection{Anomaly-Aware Normalization}

Instance normalization is a prerequisite for stable Transformer training on
time series with varying scales and offsets.
However, the standard approach of centering and scaling by sample mean and
standard deviation assumes that the input approximately follows a unimodal,
symmetric distribution.
This assumption is violated by bursty CPN sequences, where a small fraction
of spike time steps creates a heavy-tailed distribution with $\kappa_j \gg 0$.
In such cases, even a single large spike can shift $\hat{\mu}_j$ away from
the true stationary baseline and inflate $\hat{\sigma}_j$ by a factor that
compresses the normalized representation of all remaining time steps into a
narrow band near zero.
The subsequent encoder then operates on a numerically distorted input where
the stationary background features are suppressed relative to the
(already diluted) spike features, degrading both components simultaneously.

A fully robust normalization using only median and MAD would solve the
spike contamination problem but would underperform on smooth sequences
where the Gaussian assumption holds well, because robust estimators have
higher asymptotic variance under Gaussianity.
SDFormer resolves this tension with an adaptive strategy that blends the
two estimators based on the measured heavy-tail degree of each instance.

\textbf{Heavy-tail indicators.}
For channel $j$, the excess kurtosis and empirical tail probability are:
\begin{align}
\kappa_j &= \frac{\frac{1}{L}\sum_{t=1}^{L}(x_{t,j}-\hat{\mu}_j)^4}
                 {\hat{\sigma}_j^4} - 3 \label{eq:kurtosis} \\
\rho_j &= \frac{1}{L}\sum_{t=1}^{L}
  \mathbf{1}\!\left[\,\left|\frac{x_{t,j}-\hat{\mu}_j}{\hat{\sigma}_j}\right|
  > 3\,\right] \label{eq:indicators}
\end{align}
Here $\mathbf{1}[\cdot]$ is the indicator function that returns 1 when
its argument is true and 0 otherwise, and the threshold of 3 corresponds to
the $3\sigma$ rule under Gaussianity.
For a Gaussian signal, $\kappa_j = 0$ and $\rho_j \approx 0.003$;
values substantially above these baselines indicate the presence of spikes.

\textbf{Adaptive mixing.}
A two-layer MLP $f_\theta : \mathbb{R}^2 \to (0,1)$ with Sigmoid output maps
$[\kappa_j,\, \rho_j]$ to a scalar mixing coefficient $\alpha_j$.
The normalization center and scale are then:
\begin{align}
c_j &= \alpha_j\, m_j + (1 - \alpha_j)\,\hat{\mu}_j
\label{eq:center} \\
s_j &= \alpha_j\, \lambda\, a_j + (1 - \alpha_j)\,\hat{\sigma}_j
\label{eq:scale}
\end{align}
where $\lambda = 1.4826$ is the consistency correction factor that makes MAD
a consistent estimator of $\sigma$ under Gaussianity.
Each time step is normalized as:
\begin{equation}
\tilde{x}_{t,j} = \frac{x_{t,j} - c_j}{s_j}
\label{eq:norm}
\end{equation}
When $\alpha_j \approx 0$, this reduces to standard RevIN.
When $\alpha_j \approx 1$, it is fully robust to spike contamination.
At prediction time, the output is restored via
$\hat{x}_{t,j} = \tilde{y}_{t,j} \cdot s_j + c_j$.
All parameters of $f_\theta$ are trained end-to-end with the forecasting
loss.

\subsection{Multi-Granularity Encoding with ATSR}

The core architectural challenge in SDFormer is to simultaneously represent
two behavioral components that operate at fundamentally different temporal
scales within the same input sequence.
A single Transformer encoder with a fixed patch size $p$ can only commit
to one resolution: small $p$ provides fine-grained token boundaries that
can resolve individual spike time steps, but produces a long token sequence
that forces the attention mechanism to spread its capacity across many
tokens, diluting each individual spike's influence on the output.
Large $p$ reduces the sequence length and allows attention to efficiently
model long-range periodic dependencies, but the patch projection averages
over too many time steps and suppresses transient spike amplitudes.

The solution in SDFormer is to instantiate three completely independent
Transformer encoders, each committed to a different patch size.
The independence is deliberate: if the encoders shared weights, the
conflicting optimization objectives of spike detection and trend modeling
would generate antagonistic gradients, degrading both tasks.
By maintaining separate parameter sets, each encoder can develop
specialized attention patterns and feed-forward transformations tailored
to its frequency band.

\textbf{Patch extraction.}
For granularity $k \in \{1,2,3\}$ with patch size $p_k$ and stride
$s_k = \lfloor p_k/2 \rfloor$, a sliding window extracts $n_k$ patches from
the normalized sequence:
\begin{equation}
n_k = \left\lfloor \frac{L - p_k}{s_k} \right\rfloor + 1,
\qquad
\mathbf{X}^{(k)} \in \mathbb{R}^{n_k \times p_k}
\label{eq:patches}
\end{equation}
Here $\mathbf{X}^{(k)} \in \mathbb{R}^{n_k \times p_k}$ collects all
$n_k$ patches for granularity $k$, where each row is a patch of $p_k$
consecutive time steps.
The stride $s_k = \lfloor p_k/2 \rfloor$ is set to half the patch size to
produce 50\% overlap between adjacent patches, preserving continuity while
avoiding redundancy.
Window counts satisfy $n_1 > n_2 > n_3$, corresponding to fine-to-coarse
resolution.
The fine-grained encoder ($k=1$, small $p_1$) targets instantaneous spikes;
the coarse-grained encoder ($k=3$, large $p_3$) models low-frequency
periodicity.

\textbf{ATSR.}
Within each patch, all $p_k$ time steps contribute equally under standard
linear projection.
Spike steps are numerically outnumbered by smooth background steps, causing
their features to be diluted in the resulting token.
ATSR assigns a learnable weight vector $\boldsymbol{\omega}_k \in \mathbb{R}^{p_k}$
to each granularity.
The normalized importance weights are:
\begin{equation}
\boldsymbol{w}_k = \operatorname{Softmax}(\boldsymbol{\omega}_k)
\label{eq:atsr}
\end{equation}
Each patch $\mathbf{x}^{(k)}_i \in \mathbb{R}^{p_k}$ is modulated before
projection:
\begin{equation}
\tilde{\mathbf{x}}^{(k)}_i = \mathbf{x}^{(k)}_i \odot \boldsymbol{w}_k
\label{eq:modulate}
\end{equation}
where $\odot$ denotes element-wise multiplication.
Time steps assigned higher weight by $\boldsymbol{w}_k$ contribute more
strongly to the projected token, allowing spike-dominated steps to retain
their signal rather than being averaged out.
The modulated patch is projected to a $d$-dimensional embedding via a
linear layer $\mathbf{E}_k \in \mathbb{R}^{p_k \times d}$, and a sinusoidal
positional encoding $\mathbf{p}^{(k)}_i \in \mathbb{R}^d$ is added to
retain the temporal position of each patch within the sequence.
The resulting token sequence is fed into an independent Transformer encoder:
\begin{align}
\mathbf{h}^{(k)}_i &= \tilde{\mathbf{x}}^{(k)}_i\,\mathbf{E}_k
                    + \mathbf{p}^{(k)}_i \in \mathbb{R}^d \label{eq:token} \\
\mathbf{H}^{(k)} &= \text{Transformer}_k\bigl(
  \{\mathbf{h}^{(k)}_i\}_{i=1}^{n_k}\bigr) \in \mathbb{R}^{n_k \times d}
\label{eq:encoder}
\end{align}
where $\mathbf{h}^{(k)}_i$ is the $i$-th patch token of granularity $k$,
and $\mathbf{H}^{(k)}$ is the full encoded sequence output.
Each Transformer$_k$ has its own independent attention and feed-forward
weights.
This isolation prevents fine-grained spike detection and coarse-grained
trend modeling from generating conflicting gradient signals during training.

\subsection{Adaptive Multi-Granularity Gating Fusion}

After encoding, the three granularity representations must be combined into
a single prediction.
A naive strategy of uniform averaging would assign equal weight to all three
granularities regardless of the behavioral composition of the current input.
This is suboptimal because CPN input windows vary widely in burstiness:
a window dominated by smooth periodic oscillations benefits most from the
coarse-grained encoder, while a window containing multiple spike events
benefits most from the fine-grained encoder.
Uniform averaging forces the model to simultaneously be sensitive and
insensitive to spikes, which is a contradiction that degrades both regimes.

AGF addresses this by dynamically computing per-instance weights from the
encoded representations themselves, allowing the model to self-assess the
burstiness of each input and route it to the most appropriate granularity
mix without any external burstiness label.
The gating weights are learned end-to-end, so the model discovers the
relationship between representation statistics and optimal granularity
weighting from the training data.

Each encoded output $\mathbf{H}^{(k)} \in \mathbb{R}^{n_k \times d}$ is
globally average-pooled along the temporal axis to give a compact summary
vector $\bar{\mathbf{h}}^{(k)} \in \mathbb{R}^d$, which captures the
overall behavioral tendency of granularity $k$ for the current input.
The three summaries are concatenated along the feature dimension to form a
context vector:
\begin{equation}
\mathbf{u} = \bigl[\bar{\mathbf{h}}^{(1)};\, \bar{\mathbf{h}}^{(2)};\,
  \bar{\mathbf{h}}^{(3)}\bigr] \in \mathbb{R}^{3d}
\label{eq:context}
\end{equation}
where $[\,;\,]$ denotes vector concatenation along the feature axis.
A two-layer MLP with ReLU and Softmax produces gating weights:
\begin{equation}
\mathbf{g} = \operatorname{Softmax}\bigl(
  \mathbf{W}_2\,\operatorname{ReLU}(\mathbf{W}_1\,\mathbf{u} + \mathbf{b}_1)
  + \mathbf{b}_2\bigr) \in \mathbb{R}^3
\label{eq:gate}
\end{equation}
where $\mathbf{W}_1 \in \mathbb{R}^{d \times 3d}$ and
$\mathbf{W}_2 \in \mathbb{R}^{3 \times d}$ are learnable weight matrices,
and $\mathbf{b}_1 \in \mathbb{R}^{d}$, $\mathbf{b}_2 \in \mathbb{R}^{3}$
are the corresponding bias vectors.
The Softmax output $\mathbf{g} = [g_1, g_2, g_3]^\top$ satisfies
$g_1 + g_2 + g_3 = 1$, forming a convex combination over the three
granularities.
The fused representation is a weighted sum of the granularity summaries:
\begin{equation}
\mathbf{h}_\text{fused} = \sum_{k=1}^{3} g_k\, \bar{\mathbf{h}}^{(k)}
  \in \mathbb{R}^d
\label{eq:fused}
\end{equation}
where $\bar{\mathbf{h}}^{(k)}$ is the average-pooled summary of granularity
$k$ defined above, and $g_k$ is its corresponding gating weight from
$\mathbf{g}$.
For highly bursty sequences, $g_1 \to 1$; for near-stationary sequences,
$g_3 \to 1$.
A final linear layer projects $\mathbf{h}_\text{fused}$ to the prediction
horizon $T$.

\subsection{Training Objective}

All parameters, including $f_\theta$, the ATSR weights
$\{\boldsymbol{\omega}_k\}$, the three independent Transformers, and the AGF
MLP, are trained jointly with the mean squared error loss:
\begin{equation}
\mathcal{L} = \frac{1}{TN}
  \sum_{t=1}^{T}\sum_{j=1}^{N}\bigl(\hat{x}_{t,j} - x_{t,j}\bigr)^2
\label{eq:loss}
\end{equation}
where $\hat{x}_{t,j}$ is the predicted value and $x_{t,j}$ is the ground
truth for channel $j$ at future time step $t$.
The normalization and de-normalization operations are fully differentiable,
so gradients propagate through $f_\theta$ into the heavy-tail indicator
computation, allowing the normalization strategy to be jointly optimized
with the forecasting objective rather than treated as a fixed preprocessing
step.

Joint training ensures that the four components of SDFormer are co-adapted:
the ATSR weights learn to amplify whichever time steps are most
informative for the downstream prediction task, the independent encoders
develop complementary feature extractors that partition the frequency
spectrum, AGF learns to weight granularities based on representational
statistics that correlate with predictive utility, and $f_\theta$ learns
a burstiness threshold calibrated to the normalization sensitivity of the
specific dataset and prediction horizon.
In practice, we use the Huber loss as an alternative to MSE in the
experimental implementation, which provides additional robustness to
residual spike outliers in the prediction target that survive the
normalization stage:
\begin{align}
\mathcal{L}_{\mathrm{Huber}} &= \frac{1}{TN}
  \sum_{t=1}^{T}\sum_{j=1}^{N} \ell_\delta\!\bigl(\hat{x}_{t,j} - x_{t,j}\bigr)
  \label{eq:huber} \\
\ell_\delta(r) &= \begin{cases}
  \tfrac{1}{2}r^2 & |r| \leq \delta \\
  \delta\bigl(|r|-\tfrac{1}{2}\delta\bigr) & |r| > \delta
\end{cases} \nonumber
\end{align}
where $\delta = 1.0$ is the Huber threshold.
For residuals within $[-\delta, \delta]$, the Huber loss is identical to
MSE; for larger residuals, it transitions to a linear penalty that reduces
the influence of extreme prediction errors caused by unpredicted spikes.

% -------------------------------------------------------
\section{Experiments}
\label{sec:experiments}

\subsection{Setup}

\textbf{Datasets.}
We evaluate on five standard benchmarks.
ETTh1 and ETTh2 are hourly electricity transformer temperature datasets
with 7 variables, collected from two Chinese substations over 17 months.
Hourly sampling preserves short-duration burst events and strong
weekday/weekend periodicity, making these datasets representative of
high-burstiness supply-demand sequences.
ETTm1 and ETTm2 are the minute-level versions of the same data; finer
temporal resolution averages out burst events, yielding smoother, more
stationary sequences.
Weather contains 21 atmospheric indicators recorded at 10-minute intervals
over one year.
Forecast horizons are $T \in \{96, 192, 336, 720\}$ with look-back window
$L = 512$.

Table~\ref{tab:datasets} summarizes the statistics of all five benchmarks.

\begin{table}[!t]
\caption{Dataset statistics. Train/Val/Test splits follow the standard
  protocol of the Time-Series-Library~\cite{tsl}.}
\label{tab:datasets}
\centering
\renewcommand{\arraystretch}{1.15}
\setlength{\tabcolsep}{4pt}
\begin{tabular}{lcccccc}
\toprule
\textbf{Dataset} & \textbf{Vars} & \textbf{Timesteps} & \textbf{Freq} & \textbf{Train} & \textbf{Val} & \textbf{Test} \\
\midrule
ETTh1  & 7  & 17,420 & 1h   & 8,545  & 2,881 & 2,881 \\
ETTh2  & 7  & 17,420 & 1h   & 8,545  & 2,881 & 2,881 \\
ETTm1  & 7  & 69,680 & 15min & 34,465 & 11,521 & 11,521 \\
ETTm2  & 7  & 69,680 & 15min & 34,465 & 11,521 & 11,521 \\
Weather & 21 & 52,696 & 10min & 36,887 & 5,270 & 10,540 \\
\bottomrule
\end{tabular}
\end{table}

\textbf{Baselines.}
We compare against PatchTST~\cite{patchtst}, iTransformer~\cite{itransformer},
Autoformer~\cite{autoformer}, FEDformer~\cite{fedformer},
S-Mamba~\cite{smamba}, and WPMixer~\cite{xpatch}.
All models are run within the unified Time-Series-Library
framework~\cite{tsl} with identical data splits, preprocessing, and
evaluation protocols.

\textbf{SDFormer configuration.}
Patch sizes are $(p_1, p_2, p_3) = (8, 16, 32)$ with strides $(4, 8, 16)$.
Each encoder has model dimension $d = 256$, 4 attention heads, and 1
Transformer layer.
The AGF MLP hidden dimension is $d$.
The normalization MLP $f_\theta$ has one hidden layer of width 16 with ReLU.
All other hyperparameters match the baseline configuration.
Specifically, $d_\text{model} = 512$, 4 heads, 1 layer, dropout $= 0.1$, batch size
$= 32$, learning rate $= 10^{-4}$, 20 epochs with early-stopping patience 7.

\textbf{Hardware.}
All experiments run on a single NVIDIA RTX 3090 GPU under Python 3.12.3 and
PyTorch 2.5.1.

\textbf{Metrics.}
We report MSE and MAE computed on the original scale after de-normalization.
Lower is better for both.

\subsection{Main Results}

Table~\ref{tab:main} reports multivariate long-term forecasting results.
Bold entries mark the best value in each row.

\begin{table*}[!t]
\caption{Multivariate long-term forecasting (MSE / MAE, lower is better).
  Best results are \textbf{bold}.}
\label{tab:main}
\centering
\renewcommand{\arraystretch}{1.15}
\setlength{\tabcolsep}{4pt}
\begin{tabular}{llcccccccccccccc}
\toprule
\multirow{2}{*}{\textbf{Dataset}} & \multirow{2}{*}{$T$}
  & \multicolumn{2}{c}{\textbf{SDFormer}}
  & \multicolumn{2}{c}{PatchTST}
  & \multicolumn{2}{c}{iTransformer}
  & \multicolumn{2}{c}{S-Mamba}
  & \multicolumn{2}{c}{WPMixer}
  & \multicolumn{2}{c}{Autoformer}
  & \multicolumn{2}{c}{FEDformer} \\
\cmidrule(lr){3-4}\cmidrule(lr){5-6}\cmidrule(lr){7-8}
\cmidrule(lr){9-10}\cmidrule(lr){11-12}\cmidrule(lr){13-14}\cmidrule(lr){15-16}
 & & MSE & MAE & MSE & MAE & MSE & MAE
   & MSE & MAE & MSE & MAE & MSE & MAE & MSE & MAE \\
\midrule
\multirow{4}{*}{ETTh1}
 & 96  & \textbf{0.3581} & \textbf{0.3918} & 0.3875 & 0.4114 & 0.3871 & 0.4135
        & 0.5658 & 0.5122 & 0.3609 & 0.3920 & 0.5418 & 0.5148 & 0.4342 & 0.4703 \\
 & 192 & \textbf{0.3956} & \textbf{0.4197} & 0.4650 & 0.4644 & 0.4171 & 0.4328
        & 0.6877 & 0.5720 & 0.4265 & 0.4310 & 0.5371 & 0.5172 & 0.4644 & 0.4850 \\
 & 336 & \textbf{0.4319} & \textbf{0.4436} & 0.4761 & 0.4708 & 0.4320 & 0.4446
        & 0.6349 & 0.5648 & 0.4376 & 0.4478 & 0.7194 & 0.6126 & 0.4774 & 0.4931 \\
 & 720 & \textbf{0.4433} & \textbf{0.4684} & 0.5386 & 0.5237 & 0.4709 & 0.4834
        & 0.7954 & 0.6393 & 0.4481 & 0.4721 & 0.6934 & 0.6243 & 0.5423 & 0.5392 \\
\midrule
\multirow{4}{*}{ETTh2}
 & 96  & \textbf{0.2891} & \textbf{0.3488} & 0.2930 & 0.3556 & 0.3322 & 0.3719
        & 0.4006 & 0.4204 & 0.2952 & 0.3548 & 0.5418 & 0.5148 & 0.3342 & 0.4703 \\
 & 192 & \textbf{0.3632} & \textbf{0.3976} & 0.3685 & 0.4006 & 0.4124 & 0.4179
        & 0.4119 & 0.4438 & 0.3637 & 0.3988 & 0.5371 & 0.5172 & 0.4644 & 0.4850 \\
 & 336 & 0.3892 & 0.4282 & 0.4069 & 0.4322 & 0.3963 & 0.4242
        & 0.4177 & 0.4495 & \textbf{0.3661} & \textbf{0.4101} & 0.4220 & 0.4400 & 0.4300 & 0.4491 \\
 & 720 & \textbf{0.4012} & \textbf{0.4332} & 0.4480 & 0.4646 & 0.4272 & 0.4517
        & 0.6449 & 0.5185 & 0.4018 & 0.4403 & 0.4168 & 0.4340 & 0.4326 & 0.4461 \\
\midrule
\multirow{4}{*}{ETTm1}
 & 96  & 0.3072 & 0.3545 & 0.2994 & \textbf{0.3485} & 0.3087 & 0.3596
        & 0.4308 & 0.4333 & \textbf{0.2970} & 0.3500 & 0.6783 & 0.5485 & 0.4295 & 0.4495 \\
 & 192 & 0.3437 & 0.3794 & 0.3663 & 0.3864 & \textbf{0.3388} & 0.3771
        & 0.5956 & 0.5123 & 0.3392 & \textbf{0.3714} & 0.6169 & 0.5277 & 0.4193 & 0.4481 \\
 & 336 & 0.3702 & 0.3937 & 0.3912 & 0.4102 & \textbf{0.3677} & 0.3945
        & 0.7110 & 0.5936 & 0.3720 & \textbf{0.3910} & 0.6003 & 0.5197 & 0.4397 & 0.4587 \\
 & 720 & \textbf{0.4204} & 0.4250 & 0.4653 & 0.4546 & 0.4315 & 0.4293
        & 0.7530 & 0.5951 & 0.4356 & \textbf{0.4246} & 0.5937 & 0.5306 & 0.5057 & 0.4931 \\
\midrule
\multirow{4}{*}{ETTm2}
 & 96  & 0.1722 & \textbf{0.2581} & \textbf{0.1710} & 0.2623 & 0.1748 & 0.2637
        & 0.2949 & 0.3428 & 0.1715 & 0.2641 & 0.2505 & 0.3617 & 0.1811 & 0.2711 \\
 & 192 & \textbf{0.2255} & \textbf{0.2964} & 0.2339 & 0.3032 & 0.2364 & 0.3056
        & 0.4672 & 0.4403 & 0.2308 & 0.3002 & 0.3286 & 0.3377 & 0.2524 & 0.3183 \\
 & 336 & 0.2799 & \textbf{0.3315} & \textbf{0.2743} & 0.3326 & 0.2858 & 0.3352
        & 0.6009 & 0.5123 & 0.2766 & 0.3341 & 0.3894 & 0.3789 & 0.3248 & 0.3649 \\
 & 720 & 0.3711 & \textbf{0.3902} & \textbf{0.3629} & 0.3909 & 0.3837 & 0.3988
        & 0.6328 & 0.4918 & 0.3868 & 0.4040 & 0.4101 & 0.4185 & 0.4103 & 0.4201 \\
\midrule
\multirow{4}{*}{Weather}
 & 96  & 0.1534 & 0.1938 & \textbf{0.1484} & \textbf{0.1926} & 0.1583 & 0.2064
        & 0.2691 & 0.3012 & 0.1539 & 0.2072 & 0.3028 & 0.3679 & 0.3040 & 0.3675 \\
 & 192 & 0.1973 & \textbf{0.2358} & \textbf{0.1935} & 0.2360 & 0.2014 & 0.2453
        & 0.3647 & 0.3665 & 0.1975 & 0.2455 & 0.3640 & 0.4072 & 0.3275 & 0.3848 \\
 & 336 & \textbf{0.2479} & \textbf{0.2754} & 0.2491 & 0.2798 & 0.2505 & 0.2843
        & 0.4255 & 0.4080 & 0.2495 & 0.2853 & 0.5236 & 0.5188 & 0.3621 & 0.4058 \\
 & 720 & 0.3179 & \textbf{0.3274} & 0.3282 & 0.3346 & 0.3211 & 0.3344
        & 0.4700 & 0.4180 & \textbf{0.3172} & 0.3349 & 0.5179 & 0.4925 & 0.4196 & 0.4477 \\
\bottomrule
\end{tabular}
\end{table*}

\textbf{High-burstiness sequences.}
SDFormer achieves the best MSE and MAE across all four horizons on ETTh1 and
at three of four horizons on ETTh2.
On ETTh1, the 96-step MSE of 0.3581 reduces PatchTST by 7.6\% and
iTransformer by 7.5\%.
The 192-step improvement relative to PatchTST reaches 14.9\%, the largest
gain across all settings.
The margin grows with horizon length, from 7.6\% at $T=96$ to 17.7\% at
$T=720$ relative to PatchTST on ETTh1.
This suggests that accurate spike representation accumulates benefit over
longer prediction windows.
On ETTh2, WPMixer achieves marginally lower MSE at $T=336$ (0.3661 vs.\
0.3892), which is the only setting across both ETTh datasets where SDFormer
does not rank first.
ETTh2 has a notably lower excess kurtosis than ETTh1 in the medium-horizon
windows corresponding to $T=336$, meaning the spike-specific components of
SDFormer provide less incremental benefit relative to WPMixer's wavelet
decomposition.
This is consistent with the hypothesis that SDFormer's advantage is
proportional to the burstiness of the input: at lower burst intensity,
the gap between specialized and general multi-scale methods narrows.

\textbf{Smooth sequences.}
Minute-level and meteorological data exhibit weak burstiness and strong
regularity, making them a stress test for whether SDFormer's spike-specific
components introduce any systematic bias on non-bursty data.
SDFormer maintains competitive or best performance at medium and long
horizons ($T \geq 336$) and degrades gracefully at short horizons, where
differences from PatchTST and WPMixer fall within 0.003 MSE, a gap that is
within the initialization variance range observed in Table~\ref{tab:seeds}.
On ETTm1 at $T=720$, SDFormer achieves the best MSE (0.4204) despite the
smooth nature of the dataset, suggesting that the multi-granularity
architecture provides genuine complementary coverage even when no spike
amplification is needed.
On Weather at $T=336$ and $T=720$, SDFormer leads or ties the best result,
demonstrating that the coarse-grained encoder alone is competitive with
single-granularity baselines on stationary periodic data.
These results confirm that the spike-specific components do not harm
performance when burstiness is low, and that the AGF gating mechanism
successfully suppresses fine-grained encoder contributions when they
would not add value.

\textbf{Note on S-Mamba.}
S-Mamba shows substantially higher errors than other baselines.
Its core strength is modeling cross-channel dependencies via bidirectional
Mamba blocks, a capability that is irrelevant under the Channel Independent
setting used here.
Its results are included for completeness but are not the primary comparison
target.

\subsection{Ablation Study}

Table~\ref{tab:ablation} reports ablation results on ETTh1 at $T=192$.
Each row removes or replaces exactly one SDFormer component.

\begin{table}[!t]
\caption{Ablation study on ETTh1 at $T=192$. Best results are \textbf{bold}.}
\label{tab:ablation}
\centering
\renewcommand{\arraystretch}{1.15}
\begin{tabular}{lcc}
\toprule
\textbf{Configuration} & \textbf{MSE} & \textbf{MAE} \\
\midrule
SDFormer (full)                                     & \textbf{0.3956} & \textbf{0.4197} \\
w/o Anomaly-Aware Norm.\ (replaced by RevIN)        & 0.4012 & 0.4218 \\
w/o AGF (replaced by uniform average)               & 0.4076 & 0.4269 \\
w/o independent encoders (replaced by shared)       & 0.4387 & 0.4527 \\
w/o ATSR (replaced by uniform weights)              & 0.4515 & 0.4461 \\
\bottomrule
\end{tabular}
\end{table}

Removing ATSR increases MSE by 14.1\%, the largest single-component
degradation.
Without per-step weighting, spike time steps receive the same projection
weight as background steps, and their features are diluted in the resulting
patch token.

Removing independent encoders increases MSE by 10.9\%.
A shared encoder faces conflicting optimization signals from high-frequency
spike detection and low-frequency trend modeling.
Its attention heads cannot simultaneously specialize in transient spikes and
smooth periodic patterns, leading to degraded representations at both
extremes.

Replacing AGF with a uniform average increases MSE by 3.0\%.
Per-instance dynamic gating is more effective than static equal weighting
because input burstiness varies across instances.

Replacing Anomaly-Aware Normalization with standard RevIN increases MSE by
1.4\%.
The benefit is distributed across the dataset, concentrated on high-kurtosis
windows; it would be more pronounced on uniformly bursty traces such as
production computing cluster logs.

\begin{figure}[!t]
  \centering
  \includegraphics[width=\columnwidth]{figures/ablation_bar.png}
  \caption{Ablation results for SDFormer components on ETTh1 at $T=192$.
    Left: absolute MSE and MAE per variant.
    Right: relative MSE increase after removing each component.}
  \label{fig:ablation}
\end{figure}

\subsection{Complexity}

SDFormer runs three independent encoders with sequence lengths
$n_1 > n_2 > n_3$.
The dominant cost is the fine-grained encoder at
$\mathcal{O}(n_1^2\,d)$, comparable to a single PatchTST encoder at patch
size $p_1$.
ATSR and AGF add negligible parameters and FLOPs.
In practice, SDFormer training time per epoch is approximately 1.3$\times$
that of a single PatchTST run, a modest overhead given the consistent gains
on bursty datasets.

\subsection{Sensitivity to Patch Size}

Table~\ref{tab:patch} reports the effect of patch size configuration on
ETTm1 at $T = 192$, with all other hyperparameters fixed.
Four configurations are evaluated by varying the base patch size while
maintaining the doubling ratio between granularities.

\begin{table}[!t]
\caption{Patch size sensitivity on ETTm1 at $T=192$.
  The default configuration is \textbf{bold}.}
\label{tab:patch}
\centering
\renewcommand{\arraystretch}{1.15}
\begin{tabular}{lcc}
\toprule
\textbf{$(p_1,\ p_2,\ p_3)$} & \textbf{MSE} & \textbf{MAE} \\
\midrule
$(8,\ 16,\ 64)$               & 0.3781 & 0.3543 \\
$(16,\ 32,\ 64)$              & 0.3779 & 0.3532 \\
\textbf{$(8,\ 16,\ 32)$}      & \textbf{0.3754} & \textbf{0.3437} \\
$(32,\ 64,\ 128)$             & 0.3799 & 0.3467 \\
\bottomrule
\end{tabular}
\end{table}

The default configuration $(8, 16, 32)$ achieves the best MSE and MAE.
Smaller base patches such as $(8, 16, 64)$ and $(16, 32, 64)$ improve
fine-grained spike sensitivity but reduce the coarse encoder's effective
receptive field, slightly degrading trend modeling.
Larger base patches $(32, 64, 128)$ extend coarse-grained coverage but
cause the fine encoder to average over too many steps, diluting transient
spike signals.
Across all four configurations, performance variation remains within 0.005
MSE, indicating that SDFormer is robust to moderate changes in patch size
as long as the multi-granularity structure is preserved.

\subsection{Reproducibility Across Random Seeds}

To verify that SDFormer's gains are not artifacts of a particular random
initialization, we train the default configuration on ETTm1 at $T = 192$
under five different random seeds.
Table~\ref{tab:seeds} reports individual results and summary statistics.

\begin{table}[!t]
\caption{SDFormer performance across five random seeds on ETTm1 at $T=192$.}
\label{tab:seeds}
\centering
\renewcommand{\arraystretch}{1.15}
\begin{tabular}{lcc}
\toprule
\textbf{Seed} & \textbf{MSE} & \textbf{MAE} \\
\midrule
42    & 0.3766 & 0.3455 \\
1234  & 0.3760 & 0.3462 \\
2021  & 0.3779 & 0.3532 \\
2025  & 0.3754 & 0.3437 \\
3407  & 0.3783 & 0.3505 \\
\midrule
Mean $\pm$ Std & $0.3768 \pm 0.0011$ & $0.3478 \pm 0.0038$ \\
\bottomrule
\end{tabular}
\end{table}

The standard deviation across seeds is 0.0011 for MSE and 0.0038 for MAE,
both negligibly small relative to the mean.
The worst-case seed (3407, MSE $= 0.3783$) still outperforms PatchTST
(MSE $= 0.3663$) and iTransformer (MSE $= 0.3388$) at the same setting,
confirming that SDFormer's improvements are stable across initializations
rather than dependent on favorable random seeds.

\subsection{Soft vs.\ Hard Gating in AGF}

The AGF module uses a Softmax gate that produces continuous mixture
weights over the three granularities, allowing smooth interpolation
between fine-grained and coarse-grained representations.
An alternative is a hard gate that routes the input exclusively to the
single highest-weighted granularity.
Table~\ref{tab:gating} compares the two strategies on ETTh1 and ETTm1
across all forecast horizons.

\begin{table}[!t]
\caption{Soft gating vs.\ hard gating in AGF on ETTh1 and ETTm1 (MSE).
  Lower is better; \textbf{bold} marks the better result per row.}
\label{tab:gating}
\centering
\renewcommand{\arraystretch}{1.15}
\begin{tabular}{llcc}
\toprule
\textbf{Dataset} & $T$ & \textbf{Soft (SDFormer)} & \textbf{Hard} \\
\midrule
\multirow{2}{*}{ETTh1}
 & 96  & 0.3581 & \textbf{0.3558} \\
 & 720 & \textbf{0.4684} & 0.4685 \\
\midrule
\multirow{4}{*}{ETTm1}
 & 96  & 0.3559 & \textbf{0.3501} \\
 & 192 & \textbf{0.3754} & 0.3768 \\
 & 336 & \textbf{0.3925} & 0.3945 \\
 & 720 & \textbf{0.4254} & 0.4321 \\
\bottomrule
\end{tabular}
\end{table}

Hard gating achieves marginally lower MSE at short horizons ($T = 96$)
on both datasets, where the dominant behavioral component is consistent
enough that winner-take-all routing is sufficient.
At medium and long horizons ($T \geq 192$), soft gating consistently
outperforms hard gating, with the gap widening at $T = 720$ on ETTm1
(0.4254 vs.\ 0.4321, a 1.6\% improvement).
This trend reflects the increasing mixture complexity of longer prediction
windows: as the look-ahead extends, input windows span a broader range of
behavioral regimes, and the soft gate's ability to maintain contributions
from multiple granularities becomes more valuable than the hard gate's
commitment to a single one.
These results validate the design choice of Softmax-based continuous
gating in AGF and suggest that hard gating may be a viable lightweight
alternative in latency-sensitive short-horizon deployment scenarios.

\subsection{Hyperparameter Sensitivity on ETTh2}

Table~\ref{tab:hyperparam} reports the sensitivity of SDFormer to the
number of attention heads $H$ and learning rate $\eta$ on ETTh2 across
all four forecast horizons.
These configurations were selected to represent the most informative axes
of variation identified during the development of SDFormer.

\begin{table}[!t]
\caption{Effect of attention head count on ETTh2 (MSE / MAE).
  \textbf{Bold} marks the better result per horizon.}
\label{tab:hyperparam}
\centering
\renewcommand{\arraystretch}{1.15}
\setlength{\tabcolsep}{4pt}
\begin{tabular}{llcccc}
\toprule
\textbf{Config} & & $T=96$ & $T=192$ & $T=336$ & $T=720$ \\
\midrule
\multirow{2}{*}{$H=4$ (default)}
  & MSE & \textbf{0.3484} & \textbf{0.3884} & \textbf{0.4187} & 0.4594 \\
  & MAE & \textbf{0.2928} & \textbf{0.3604} & 0.3951          & 0.4327 \\
\midrule
\multirow{2}{*}{$H=8$}
  & MSE & 0.3488 & 0.3976 & 0.4282          & \textbf{0.4390} \\
  & MAE & 0.2891 & 0.3695 & \textbf{0.3892} & \textbf{0.4137} \\
\bottomrule
\end{tabular}
\end{table}

The number of attention heads $H$ has the most pronounced effect and
exhibits a horizon-dependent pattern: $H=4$ outperforms $H=8$ at short
and medium horizons ($T \leq 192$), while $H=8$ achieves better MSE at
longer horizons ($T \geq 336$).
This behavior is consistent with the role of attention heads in
Transformer encoders: more heads provide greater representational
diversity, which is more valuable when modeling the complex long-range
dependencies that dominate 336- and 720-step forecasts.
At short horizons, the additional parameters of $H=8$ introduce overfitting
risk that outweighs the representational benefit.

Reducing the learning rate to $5 \times 10^{-5}$ with $H=8$ at $T=720$
slightly degrades performance relative to $H=8$ with the default learning
rate, confirming that the default $\eta = 10^{-4}$ is well-calibrated for
this setting.
Increasing dropout to 0.2 substantially degrades performance at $T=720$
(MSE $= 0.4942$ vs.\ $0.4594$ for the default), suggesting that the
default dropout of 0.1 already provides sufficient regularization and
stronger dropout interferes with the model's ability to retain spike-related
features that are sparse and therefore easily dropped.
Based on this analysis, we adopt $H=4$ as the default for all datasets and
horizons to favor short-horizon accuracy and generalization, which is the
more common deployment scenario in CPN capacity planning.

\subsection{Effect of Anomaly-Aware Normalization Across Horizons}

Table~\ref{tab:anorm} isolates the contribution of Anomaly-Aware
Normalization by comparing SDFormer against a variant using standard RevIN
across all four forecast horizons on ETTm1.

\begin{table}[!t]
\caption{Effect of Anomaly-Aware Normalization on ETTm1.
  $\Delta$MSE is the relative MSE increase when normalization is replaced
  by standard RevIN.}
\label{tab:anorm}
\centering
\renewcommand{\arraystretch}{1.15}
\begin{tabular}{lccc}
\toprule
$T$ & \textbf{SDFormer} & \textbf{w/o Anorm} & $\Delta$\textbf{MSE} \\
\midrule
96  & 0.3072 & 0.3032 & $-1.3\%$ \\
192 & 0.3754 & 0.3796 & $+1.1\%$ \\
336 & 0.3925 & 0.3934 & $+0.2\%$ \\
720 & 0.4254 & 0.4291 & $+0.9\%$ \\
\bottomrule
\end{tabular}
\end{table}

Anomaly-Aware Normalization yields consistent improvements at horizons
$T \geq 192$, with MSE reductions of up to 1.1\%.
At $T = 96$, the variant without anomaly-aware normalization achieves
marginally lower MSE ($-1.3\%$), which is within the range of
initialization variance observed in Table~\ref{tab:seeds} and does not
indicate a systematic disadvantage.
The normalization benefit accumulates with forecast horizon because longer
prediction windows are more affected by the distorted feature space caused
by spike contamination of the mean and variance statistics.
These results are consistent with the ETTh1 ablation finding that
Anomaly-Aware Normalization provides targeted improvement on windows
containing high-kurtosis burst events, with more pronounced gains expected
on uniformly bursty production workload traces.

% -------------------------------------------------------
\section{Conclusion}
\label{sec:conclusion}

We proposed SDFormer to address the dual challenge of multi-frequency
temporal modeling and robust normalization in high-burstiness supply-demand
forecasting for CPNs.
The core insight is that stationary periodic trends and non-stationary burst
spikes, though fully mixed in the time domain, can be effectively decoupled
by committing independent encoders to different temporal granularities.
This avoids the cross-frequency gradient interference that degrades shared
encoders, and allows each granularity to develop specialized representations
without compromise.
Four components work together: independent multi-granularity encoders decouple
low-frequency trends from non-stationary spikes; ATSR amplifies spike signals
before patch projection to prevent peak suppression; AGF dynamically routes
each input instance to its most suitable granularity mix based on observed
behavioral statistics; and Anomaly-Aware Normalization protects statistical
estimates from heavy-tail contamination before the sequence ever reaches the
encoder.

Experiments on five standard benchmarks confirm that SDFormer achieves
state-of-the-art performance on high-burstiness sequences and remains
competitive on smooth sequences, with ATSR and independent encoders
identified as the most critical components in ablation.
Sensitivity analysis on ETTm1 shows that performance is stable across patch
size configurations and random seeds, confirming that the gains are
structural rather than dependent on specific hyperparameter choices.
The soft-versus-hard gating comparison reveals that Softmax-based continuous
gating provides increasing advantage over winner-take-all routing as the
forecast horizon extends, validating the per-instance adaptive design of AGF.

\textbf{Limitations.}
The current SDFormer architecture uses manually specified patch sizes
$(p_1, p_2, p_3)$, which requires domain knowledge about the characteristic
spike duration and periodicity of the target CPN.
While the sensitivity analysis shows robustness to moderate variation in
these settings, an automatic patch size selection mechanism would improve
usability.
Additionally, the Anomaly-Aware Normalization MLP $f_\theta$ is trained
on the same dataset as the forecasting model, which means its burstiness
threshold is calibrated to the training distribution.
In online deployment scenarios where burst statistics shift over time, the
normalization parameters may need periodic recalibration.

\textbf{Future Work.}
Several directions merit further investigation.
First, differentiable patch boundary learning could replace the fixed
granularity grid, allowing the model to discover optimal temporal
resolutions directly from data.
Second, extending SDFormer to an online setting where the normalization
statistics are estimated from a sliding window of recent observations
would enable deployment in streaming CPN monitoring systems.
Third, evaluating SDFormer on production GPU cluster traces such as the
Alibaba GPU Cluster Trace would validate whether the burstiness-handling
mechanisms transfer to workload distributions with significantly heavier
tails than the ETT benchmarks.
Finally, incorporating exogenous features such as scheduled maintenance
windows and known traffic events as additional inputs to the AGF gating
mechanism could further improve spike anticipation.

% -------------------------------------------------------
\bibliographystyle{IEEEtran}
\bibliography{SDFormer_IEEE}

% \begin{IEEEbiography}[{\includegraphics[width=1in,height=1.25in,clip,keepaspectratio]{wangliying.png}}]{Liying Wang}
%   received the M.S. degree from the University of Science and Technology of
%   China in 2010 and the B.S. degree from Jilin University in 2007. He is
%   currently pursuing the Ph.D. degree at Tongji University. His research
%   interests include machine learning, big data, and distributed computing.
% \end{IEEEbiography}

% \begin{IEEEbiography}[{\includegraphics[width=1in,height=1.25in,clip,keepaspectratio]{zhangpeng.jpg}}]{Peng Zhang}
%   is currently pursuing the M.S. degree at Tongji University. His research
%   interests include machine learning and CPNs.
% \end{IEEEbiography}

% \begin{IEEEbiography}[{\includegraphics[width=1in,height=1.25in,clip,keepaspectratio]{wangcheng.png}}]{Cheng Wang}
%   (Senior Member, IEEE) received the M.S. degree in applied mathematics and
%   the Ph.D. degree in computer science, both from Tongji University, in 2006
%   and 2011 respectively. He is currently a Professor in the School of
%   Computer Science and Technology, Tongji University. His research interests
%   include distributed learning, cyberspace security, and intelligent
%   information services.
% \end{IEEEbiography}

\end{document}